\section{Threats to Validity}
\label{sec:threats}

The threats taxonomy follows Wohlin et al.~\cite{wohlin2012experimentation}.

\textbf{Conclusion validity.} The evaluation comprises 2{,}700
window-seed-strategy observations summarised as point estimates
(25-seed cell means); seed-level spread is available in the frozen
artifact. The pre-registered hypothesis
decisions follow the locked tolerances. The margins by which H1
passes ($17.8$~pp above the $0.80$ threshold) and H4 fails
($78$~pp below the $0.80$ threshold) are large relative to
seed-level variability, so the hypothesis decisions are unlikely
to be artefacts of seed noise.

\textbf{Internal validity.}

\textit{(i) Failure-mode distribution.} The pre-registered patch
action distribution (92/5/3) is drawn from public sysadmin
literature but is not seeded by direct measurement. Real distributions
condition on package, OS, patch type, and organisational maturity.
A different per-action distribution would shift the absolute MTTR
and rollback rate numbers; the qualitative claim (AutoHeal hits
safety bounds at high capacity) is robust to the distribution
because the safety hard-stop fires on the observed empirical rate.

\textit{(ii) Health-check fidelity.} AutoHeal's verifier
(\S\ref{sec:architecture}) approximates real post-action probes
deterministically from the action outcome. A real-world deployment
would run configuration drift checks, performance regression tests,
and dependency-graph integrity checks. The simulator's
``health\_check $==$ success'' coupling is a pessimistic baseline:
real verifiers may catch failures the action outcome alone misses,
but they may also produce false negatives that AutoHeal would
treat as silent successes. The pre-registered hard-stop is robust
to the abstraction because it triggers on the observed empirical
rate regardless of whether individual checks are accurate.

\textit{(iii) Cascading-failure detector is heuristic.} The
detector (\S\ref{sec:architecture}) flags clusters of rollbacks
sharing a CVE-id prefix. Real cascading failures correlate via
package or dependency rather than CVE id; the heuristic is a
proxy. A pre-registered alternative based on actual
dependency-graph correlations is left as future work.

\textit{(iv) Safety-bound enforcement not exercised.} The
hard-stop safety bound is \emph{detected and flagged} in the
simulation (134 window-level flags); the registered enforcement
response --- hard stop with fallback to human-in-loop for the
remaining windows --- was not exercised. The evaluated
implementation suppresses new-CVE intake for the flagged window
and continues remediation (\S\ref{sec:safety_bounds},
\S\ref{sec:deviations}). The evaluation therefore measures
detection of the safety bounds, not the coverage or throughput
consequences of enforcing them; this is a limitation of the
evaluation, stated plainly.

\textit{(v) Selection-policy coupling.} As in Papers~5--9,
HygienePrio under fixed weights drives fleet evolution for all
three strategies. AutoHeal's triage classes select a subset of
HygienePrio's top-$K$; Human-in-loop uses HygienePrio's top-30.
A closed-loop evaluation where each strategy drives its own
trajectory is the natural follow-up (Paper~9~\cite{paper9selftraj}
quantifies the bias).

\textbf{Construct validity.} Coverage is defined over EPSS$>0.5$
pairs at the moment of detection; this is the operational target
implicit in CISA BOD 22-01~\cite{cisa2021bod2201}. MTTR is
measured in windows (not days) to match the Paper~5
simulator's~\cite{paper5temporal}
clock. The Human-in-loop baseline at $K_{\text{human}} = 30$ is
calibrated to Verizon's 43-day DBIR statistic but does not model
the qualitative properties of human review (judgment calls,
business-context awareness) that real teams contribute beyond
raw capacity.

\textbf{External validity.} All claims are bounded to the synthetic
EEHDA evaluation context with real CVE attribute distributions.
External validation on real fleet telemetry remains the program's
most consequential open problem, identified throughout
Papers~3--9. AutoHeal's pre-registered safety bounds and
hypothesis-rejection stop rules are designed to be transferable
to a real-deployment evaluation, but their specific cell-mean
numbers are not.
